\scanpage{14}\mspage{12}
on peut néanmoins définir de telles flèches, comme résultant d'une \deleted{application} \inserted{flèche} (beaucoup plus riche) \deleted{de} ``composition'' de chemins entre $a'$ et $a$, et $a$ et $b$, ou entre $a$ et $b$, et $b$ et $b'$, soit en général donc l'accouplement
\[
C_{\infty}(a,b)\times C_{\infty}(b,c)\longrightarrow C_{\infty}(a,c),
\]
dont la définition directe n'est pas moins problématique, et \unsure{on|ne} peut en tous cas \uline{pas} le \deleted{faire} \inserted{donner} comme une flèche dans $\mathrm{Cat}$ \unsure{tout-même|elle-même|tout-mieux}, ou dans $\mathrm{Ind}\,\mathrm{Cat}$.

\deleted{\unsure{Pour y}} Interprétant \unsure{p.ex} \deleted{\unsure{des chemins}} \inserted{\unsure{les|ces}} $\mathrm{Ch}_{\infty}(a,b)$ en termes de chemins \uline{infinis} (des deux côtés) de $a$ à $b$, en un sens évident (mais ``essentiellement finis'' des deux côtés), on peut définir aussi
\[
\mathrm{Ch}_{\infty}^{\deleted{\unsure{+|\Sigma}}-}(a,b),\quad\mathrm{Ch}_{\infty}^{+}(a,b)
\]
comme les sous-catégories pleines formées \deleted{des} respectivement de chemins qui sont \inserted{\unsure{t.t.|tôt}} dégénérés en degrés $\geq0$ (i.e. $a_0\leftarrow a_1\to a_2\leftarrow a_3\ldots$ sont tous des identités), ce qui revient aux chemins infinis seulement à gauche, i.e. s'arrêtant à $a_0$, et symétriquement pour $\mathrm{Ch}_{\infty}^{+}(a,b)$. Même \unsure{tabac|chose|topo} pour les $\mathrm{CH}^{+},\mathrm{CH}^{-}_{\unsure{\infty}}$.
% archive sheet number: 13

